\section{Conclusion et perspectives}
\label{sec:conclusion}

\textbf{Récapitulatif.} À partir de l'unique entrée formelle $\sper = 1/2$, et sans aucun paramètre ajustable, la Théorie de la Persistance fournit, dans cet article :
\begin{itemize}[nosep,leftmargin=*]
\item un \emph{cadre catégorique minimal} : la catégorie $L^0$ et le lemme d'entropie maximale $\Lzero$ (\ref{sec:cadre}) ;
\item sept \emph{théorèmes structurants} $\Tzero$ (champ dynamique), $\Tone$ (transitions interdites), $\Ttwo$ (conservation spectrale), $\Tthree$ (transfert antidiagonal), $\Tfour$ (convergence spectrale), $\Tfive$ (point fixe unique $\mustar = 15$), $\Tsix$ (unicité complète) (\ref{sec:theoremes}) ;
\item trois \emph{axiomes de pont prouvés} $\BAthree$ (formule d'holonomie), $\BAfour$ (critère d'activation) et $\BAfive$ (produit de Pontryagin) (\ref{sec:bridge}) ;
\item une \emph{taxonomie complète des nombres premiers} : actifs $\{3, 5, 7\}$, échos $\{11, 13\}$, super-échos $\{17, 19, 23\}$, inactifs $p \geq 29$, et le canal de parité $\{2\}$ (\ref{sec:bridge:taxonomie}) ;
\item un \emph{principe de cascade arithmétique} qui agence les trois premiers actifs en une séquence ordonnée de raffinements, localement forcée par les théorèmes de pont BT$_7$ et BT$_9$ (\ref{sec:cascade}) ;
\item une \emph{dérivation complète} de la constante de structure fine, $1/\alpha_{\rm EM}^{\rm PT} = 137{,}035\,999\,083$, avec un écart à la valeur mesurée inférieur à $0{,}004$ ppb (\ref{sec:alphaem}).
\end{itemize}

\textbf{Statut épistémique.} Tous les énoncés présentés dans cet article (T0--T6 et BA0--BA5) sont, à ce jour, démontrés et indexés \textsc{[Thm]} dans la carte de nomenclature du monographe~\cite[\texttt{NOMENCLATURE\_MAP.md}]{PT_Monograph}. Le programme PT a en outre fermé en avril~2026 ses dix engagements structurels critiques (compteur \emph{C1--C10}, $10/10\;\textsc{[Thm]}$ après les promotions C2/C4/C5/C7 du 26 avril 2026). La seule étape conditionnelle qui subsiste dans l'écosystème PT et qui touche \emph{indirectement} cet article est la \emph{revue par les pairs externes}, pour laquelle ce duo d'articles constitue le matériel d'entrée.

\textbf{Article compagnon~II.} La lecture purement arithmétique présentée ici se prolonge naturellement par une \emph{lecture géométrique et spectrale} de la PT, exposée dans le second article du programme~\cite{PT_Article2_GeoSpectral}. On y trouve, en particulier :
\begin{itemize}[nosep,leftmargin=*]
\item le \emph{théorème de la courbe de persistance} : la cascade canonique~\eqref{eq:cascade_canonique} admet une unique courbe algébrique de Kontsevich--Norbury, de genre $\mustar - 1 = 14$ ;
\item le \emph{théorème du quasi-soliton de Ricci} : la métrique de Fisher--Bianchi induite par la cascade satisfait $\mathrm{Ric} + \mathrm{Hess}(f) = \lambda(\mu)\, g$ avec $\lambda(\mu) = -1/\mu^{4} + O(1/\mu^{5})$, coefficient $-1$ exact ;
\item l'\emph{identité de phase de Maslov--Higgs} : $1/8 = \sper^{2}/2 = \lambda_H = \Delta_N^{\rm Maslov}$, reliant le couplage de Higgs au résidu de Riemann--von~Mangoldt.
\end{itemize}

\textbf{Applications aval.} Le cadre développé dans cet article s'applique sans modification à un grand nombre d'observables physiques et de domaines scientifiques. Sans développement~ici, signalons les directions actives au sein du programme~PT~\cite{PT_Monograph}~:
\begin{itemize}[nosep,leftmargin=*]
\item \emph{Couplages du Modèle Standard} : $\sin^{2}\theta_W$, $\alpha_s(M_Z)$, hiérarchie des Yukawa, matrice CKM (\cite[ch.~15, ch.~18]{PT_Monograph}).
\item \emph{Chimie} : énergies de liaison, déplacements chimiques RMN, aromaticité, structure fine atomique (bibliothèque libre \texttt{PT\_CHEMISTRY}~\cite[GitHub]{PT_Monograph}).
\item \emph{Biologie quantitative} : exposants allométriques, allostérie protéique, longueurs minimales des motifs ARN.
\item \emph{Cosmologie} : tension de Hubble, structure du secteur sombre, densité de matière noire.
\item \emph{Au-delà du Modèle Standard} : taxonomie des candidats BSM, neutrinos lourds, masse du graviton, identités de Wilson.
\end{itemize}

Chacune de ces directions exploite la même chaîne $\sper = 1/2 \to \Tone\text{--}\Tsix \to \mustar = 15 \to \BAthree\text{--}\BAfive$ que l'article présent décrit, en couplant en aval les axiomes de pont aux observables physiques propres au domaine considéré. À ce titre, cet article fournit le socle commun à toutes les applications de la PT.
